\begin{tikzpicture}[font=\small]
  % ---------- array <-> tree correspondence ----------
  \begin{scope}
    \node[anchor=south, font=\small] at (4.6,3.6){二叉堆的核心技巧：\textbf{完全二叉树可以用数组隐式表示}，不需要任何指针};

    % tree
    \node[vtx] (n0) at (2.6,2.6) {$9$};
    \node[vtx] (n1) at (1.3,1.6) {$7$};
    \node[vtx] (n2) at (3.9,1.6) {$6$};
    \node[vtx] (n3) at (0.6,0.6) {$4$};
    \node[vtx] (n4) at (2.0,0.6) {$5$};
    \node[vtx] (n5) at (3.3,0.6) {$2$};
    \foreach \a/\b in {n0/n1, n0/n2, n1/n3, n1/n4, n2/n5} \draw[edge] (\a) -- (\b);

    % array
    \foreach \i/\v in {0/$9$, 1/$7$, 2/$6$, 3/$4$, 4/$5$, 5/$2$}{
      \node[cell] at (6.0+\i*0.85,1.9) {\v};
      \node[note, anchor=north] at (6.0+\i*0.85,1.5) {\i};
    }
    \node[note, anchor=south] at (8.1,2.35){数组下标};

    \node[note, anchor=west, align=left] at (6.0,0.6)
      {$\mathrm{parent}(i)=\lfloor(i-1)/2\rfloor$\\
       $\mathrm{left}(i)=2i+1$\\
       $\mathrm{right}(i)=2i+2$};

    \node[note, anchor=north, align=center] at (4.6,-0.3)
      {堆序性质：\textbf{每个结点都 $\ge$ 它的两个孩子}（最大堆）。\;
       注意这\textbf{不是}排序 —— 左右孩子之间没有任何约束。};
  \end{scope}

  % ---------- sift down ----------
  \begin{scope}[yshift=-4.0cm]
    \node[anchor=south, font=\small] at (5.0,2.4){\texttt{max\_heapify}（下沉）：把一个违反堆序的结点往下换};

    \node[bad] (a0) at (0.8,1.6) {$3$};
    \node[vtx] (a1) at (0.0,0.6) {$7$};
    \node[vtx] (a2) at (1.6,0.6) {$6$};
    \draw[edge] (a0) -- (a1); \draw[edge] (a0) -- (a2);
    \node[note, anchor=north] at (0.8,-0.1){与\textbf{较大}的孩子交换};

    \node[note] at (2.8,1.0){$\Longrightarrow$};

    \node[vtx] (b0) at (4.6,1.6) {$7$};
    \node[bad] (b1) at (3.8,0.6) {$3$};
    \node[vtx] (b2) at (5.4,0.6) {$6$};
    \draw[edge] (b0) -- (b1); \draw[edge] (b0) -- (b2);
    \node[note, anchor=north] at (4.6,-0.1){继续往下，直到就位};

    \node[note, anchor=west, align=left] at (6.6,1.2)
      {单次下沉：$O(\log n)$（树高）\\[5pt]
       \textbf{建堆}从最后一个非叶结点倒着下沉：\\
       朴素估计 $O(n\log n)$，\textbf{实际是 $O(n)$} ——\\
       因为绝大多数结点靠近叶子，下沉距离很短：\\
       $\sum_{h}\frac{n}{2^{h+1}}\cdot O(h)=O(n)$};

    \node[note, anchor=north, align=center] at (4.6,-1.2)
      {堆排序：建堆 $O(n)$ + 取 $n$ 次最大值各 $O(\log n)$ $=O(n\log n)$，且\textbf{原地}（额外空间 $O(1)$）。\\
       与归并排序同阶但不需要 $O(n)$ 辅助空间；代价是不稳定、缓存局部性差。};
  \end{scope}
\end{tikzpicture}
